% Exported from Math++. Compile with pdflatex.
\documentclass[11pt,a4paper]{article}
\usepackage[T1]{fontenc}
\usepackage[utf8]{inputenc}
\usepackage{lmodern}
\usepackage{amsmath,amssymb,amsthm,mathtools}
\usepackage{graphicx,booktabs,xcolor}
\usepackage[hidelinks]{hyperref}
\usepackage[margin=25mm]{geometry}
\theoremstyle{plain}
\newtheorem{theorem}{Theorem}
\newtheorem{lemma}[theorem]{Lemma}
\newtheorem{corollary}[theorem]{Corollary}
\newtheorem{proposition}[theorem]{Proposition}
\theoremstyle{definition}
\newtheorem{definition}[theorem]{Definition}
\newtheorem{example}[theorem]{Example}
\theoremstyle{remark}
\newtheorem*{remark}{Remark}
\setlength{\parindent}{0pt}\setlength{\parskip}{0.6em}
\begin{document}
\section*{Series}

Sums, infinite or finite.

\begin{equation*}
  \sum_{k=1}^{100} k = 5050
\end{equation*}
\begin{equation*}
  \sum_{n=1}^{\infty} \frac{1}{n^2} = \frac{\pi^{2}}{6}
\end{equation*}
\begin{equation*}
  \sum_{n=0}^{\infty} \frac{x^n}{n!} = e^{x}
\end{equation*}
\begin{equation*}
  \sum_{n=0}^{\infty} \left(\frac{1}{2}\right)^n = 2
\end{equation*}
A telescoping sum:

\begin{equation*}
  \sum_{n=1}^{\infty} \frac{1}{n (n + 1)} = 1
\end{equation*}
\end{document}
